% authority.tex -- Methodology 3.5: authority and safety as runtime assurance (drafted 12 Oct 2026 from the thesis
% plan of 11 Oct). The safety paragraph moved here from system_overview.tex (its label sec:safety comes with it);
% fig:owners (F5) replaces the old nodes and cascade figures. The rig switch states in brackets follow plan D6 and
% the M1 plan's D6(a)/(b): they are set per lab day and are not results.

\subsection{Authority and safety}
\label{sec:safety}

Every UAV in the air has exactly one controller sending it commands at every instant, and every phase has a way
down that ends with the UAV disarmed on the floor. Figure~\ref{fig:owners} shows the cascade and who owns what.
The planner owns the ring: it plans the ring's motion and every rod, and hands each carrier's share to that
carrier's tracker. A tracker owns its UAV's radio link; nothing else writes to it. A UAV that is not on the ring
is owned either by its own tracker (a freed carrier landing by itself, a newcomer being taken back) or by the
collaborator's controller (the newcomer through its pickup, transit, drop and descent to contact), and the
multiplexer in front of the newcomer's radio is the one place where ownership changes hands: it forwards the
collaborator's commands until the attach is confirmed, switches to this work's tracker on that tick, and latches.
Below the trackers each flight controller closes only its body-rate loop and arms, disarms and switches its magnet
on channels the tracker sets; the operator holds a disarm button that overrides everything.

\begin{figure}[htbp]
\centering
\begin{tikzpicture}[>=Stealth, font=\small,
    blk/.style={draw, thick, rounded corners=2pt, align=center, minimum height=0.85cm, font=\small},
    ours/.style={blk, fill=blue!6, draw=blue!50!black},
    his/.style={blk, fill=orange!10, draw=orange!60!black},
    fc/.style={blk, fill=gray!12, draw=gray!60},
    lab/.style={font=\scriptsize, align=center}]
  \node[blk]  (op)   at (5.6,1.7)   {Operator\\ \scriptsize disarm button};
  \node[ours] (fm)   at (5.6,0)     {Fleet manager\\ \scriptsize phases, faults, LAND};
  \node[ours] (plan) at (0,0)       {Load--rod planner\\ \scriptsize 10\,Hz; owns the ring};
  \node[ours] (trk)  at (0,-2.5)    {Tracker $i$ (carrier)\\ \scriptsize 50\,Hz; owns UAV $i$};
  \node[ours] (trkn) at (5.6,-2.5)  {Tracker (newcomer)\\ \scriptsize from the attach, or take-back};
  \node[his]  (tej)  at (11.3,-2.5) {Collaborator's controller\\ \scriptsize pickup, transit, drop,\\ \scriptsize descent to contact, lift};
  \node[ours] (mux)  at (5.6,-4.6)  {Mux\\ \scriptsize one owner, latched};
  \node[fc]   (bf)   at (0,-4.6)    {Flight controller $i$\\ \scriptsize Betaflight rate loop, magnet};
  \node[fc]   (bfn)  at (5.6,-6.3)  {Flight controller (newcomer)\\ \scriptsize Betaflight rate loop, magnet};
  \draw[->, thick] (op) -- node[lab, right] {commands} (fm);
  \draw[->, thick] (fm) -- node[lab, above] {phase} (plan);
  \draw[->, thick] (plan) -- node[lab, left] {reference,\\ rod pull} (trk);
  \draw[->, thick] (plan.south east) -- node[lab, sloped, above, pos=0.55] {once folded in} (trkn.north west);
  \draw[->, thick] (trk) -- node[lab, left] {throttle, rates,\\ arm, magnet} (bf);
  \draw[->, thick] (trkn) -- (mux);
  \draw[->, thick] (tej.south) -- ++(0,-1.25) -- node[lab, above] {until the attach} (mux.east);
  \draw[->, thick] (mux) -- node[lab, right] {one command stream} (bfn);
  \draw[->, thick, dashed] (trk) -- node[lab, above] {fault} (trkn);
  \draw[->, thick, dashed] (fm) -- node[lab, right] {LAND, drop, disarm} (trkn);
  \draw[->, thick, dashed] (fm.north west) -- ++(0,0.5) -| node[lab, pos=0.25, above] {arm, disarm, LAND} ([xshift=-0.5cm]plan.north west) |- (trk.west);
  \draw[<->, thick, dotted] (fm.east) -| node[lab, pos=0.25, above] {attach confirmed, give-up} (tej.north);
\end{tikzpicture}
\caption[The cascade and its owners]{The cascade and its owners. Blue: this work; orange: the collaborator's
controller; grey: the unmodified flight controllers. The planner owns the ring and the carriers' references; each
tracker owns one UAV's radio link; the newcomer's link is owned by the collaborator's controller until the attach
is confirmed and by this work's tracker from that tick, through a multiplexer that latches. Dashed: faults and the
fleet manager's commands; dotted: the confirmation and the give-up exchanged with the collaborator's stack. The
operator's disarm overrides everything.}
\label{fig:owners}
\end{figure}

\paragraph{The envelope and the fleet disarm.} Every tracker checks its own UAV against a flight envelope (UAV tilt
and ring tilt both 60$^\circ$, speed 3\,m/s) and checks that its motion-capture pose (0.25\,s) and the planner's
reference (1\,s) are fresh. A violation on a carrier disarms the whole fleet, not the one UAV: a single UAV losing
thrust on a shared rigid payload drags the others with it, and three UAVs cannot hold a ring hung for four. There
is no geofence. The argument for a fleet disarm rather than a geofenced recovery is the lab itself: the ring flies
at 1\,m inside a 6\,$\times$\,6\,$\times$\,4\,m cage, so a falling ring is a short fall onto a padded floor, where an
armed UAV flying on a wrong reference is not. The exceptions are UAVs that are not on the ring. A fault on a
newcomer that is not yet attached disarms the newcomer alone and lands the fleet, and a UAV that has left the ring
is pushed down to its own resting height and disarmed alone (Section~\ref{sec:reconfig}). LAND lands every UAV,
whatever it is doing. On the rig each UAV must confirm that its flight controller is armed before the fleet takes
off, and a UAV that disarms between ARM and TAKEOFF grounds the fleet, since three UAVs lifting a ring sized for
four would drag it.

\paragraph{The fall-backs, as runtime assurance.} The planner and the admission chain are the high-performance
controllers of a Simplex-style arrangement (Section~\ref{sec:lit_switching}); the fall-backs are simple, each has
its own monitor, and each ends with every UAV owned and down.
\begin{itemize}
    \item \textbf{Take-back.} A newcomer whose contact proves nothing, or whose owner gives up, is taken over by this
    work's tracker, moved clear of the ring and landed by itself (Section~\ref{sec:reconfig}).
    \item \textbf{Tilt cut.} A newcomer whose tilt passes its limit while it is not carrying the ring has its motors
    cut, so that a UAV caught on the rim or on a rod does not drag the ring; in the twin the fault was caught
    0.23--0.25\,s after an injected upset (exploratory), and the cut stays off on the rig
    (Section~\ref{sec:res_join_faults}).
    \item \textbf{LAND for every UAV.} LAND lands the carriers with the ring, a freed carrier where it is, and a
    newcomer wherever it is in its mission, each on its own path down.
    \item \textbf{Landing the fleet on a newcomer fault.} A fault on an unattached newcomer disarms it alone and
    lands the carriers at once, where the earlier rule disarmed all four in the air.
    \item \textbf{Dropping the ring.} A loss that leaves fewer than three carriers, or a ring tilt over 30$^\circ$,
    releases every magnet, a newcomer's included, and lands every UAV: 30$^\circ$ is where the controller gives up a
    ring it can no longer hold level, not a value to design to. It is on in every mode, the attach included, and it
    fired on the rig on 15 October after a carrier's magnet let go with no command (Section~\ref{sec:res_drop}).
    \item \textbf{The freed UAV disarmed alone.} A carrier that has left the ring lands by itself and is disarmed at
    its rest without waiting for the fleet, so that the fleet's disarm at the end of a landing never drops it.
\end{itemize}

\paragraph{What is tested where.} ``A tested way down in every phase'' holds in the twin, where every fall-back
above has been flown from the phase it guards, the attach phases on the earlier attach sequence
(Appendix~\ref{sec:app_faults}). On the rig it holds only for the
exits flown there. The drop rule has fired on the rig, after an uncommanded magnet release in a hover, and every
carrier then landed by itself and was disarmed alone; LAND and the take-back have also flown there; the tilt cut and the flight controller's own
disarm-to-land stay off on the rig until each has passed its review, so until then a take-back that sticks ends in
the operator's disarm. When a ring drop comes while the collaborator's controller flies the newcomer, the rig takes
the newcomer back and lands it (twin-tested), and the operator can ask the collaborator's planner to land it straight
down. Redistribution after a release has flown on the rig; the policy for a true fault on a carrier stays the fleet
stop and the drop rule.
